\section*{Chapitre \ref{ch:iv1:time-series} --- Titre}
\markboth{\protect\omTocMark{}{Réponses~: Séries temporelles}}{\protect\omTocMark{}{Réponses~: Séries temporelles}}

\begin{solution}{iq:iv1:time-series:1}
Stationnaires~: (a), (c), (e), (f). Non stationnaires~: (b) et (d).

La stationnarité faible\index{stationnarité faible} exige une moyenne constante, une variance constante et finie, et une autocovariance qui ne dépend que du retard.

(a) est vérifiée par définition. (c) est vérifiée, car $|\phi| < 1$ et le processus part de sa loi stationnaire. (e) est un bruit blanc, à savoir les accroissements de la marche aléatoire. (b) a une variance qui croît avec $t$. (d) a une moyenne qui varie, bien que la suppression de la tendance le rende stationnaire.

(f) est la surprise. Sa moyenne vaut 0, sa variance vaut $\cos^2t + \sin^2t = 1$ pour tout $t$, et sa covariance vaut $\cos(t - s)$, qui ne dépend que du retard. Il est donc stationnaire, bien que chaque trajectoire soit une vague parfaite. Il n'est pas ergodique~: une seule trajectoire ne révèle jamais que $A$ et $B$ pourraient différer, de sorte que la moyenne temporelle de $x_t^2$ donne $(A^2 + B^2)/2$ et non la variance 1, alors que la moyenne temporelle de $x_t$ retrouve bien la moyenne 0.

\iqwrong{Dire que (f) n'est pas stationnaire parce qu'il présente un motif visible, ou que (d) est stationnaire parce qu'il paraît stable autour de sa tendance.}

\iqfollow{Pourquoi la non-ergodicité importe-t-elle en pratique~?}{Les moyennes temporelles calculées sur une seule histoire peuvent alors ne pas estimer les valeurs d'ensemble. On n'observe jamais qu'une seule trajectoire du marché.}

\iqwhy{La stationnarité est souvent énoncée et rarement testée à partir de sa définition. La vague montre qu'un motif et la stationnarité peuvent coexister.}
\end{solution}

\begin{solution}{iq:iv1:time-series:2}
La variance vaut $1.96$, l'autocorrélation de retard 3 vaut $0.343$, la demi-vie vaut $1.94$ pas, et le $R^2$ vaut $0.7^{2h}$~: $0.49$ pour $h = 1$ et $0.24$ pour $h = 2$.

En prenant la variance des deux membres, on obtient $\gamma_0 = 0.49\gamma_0 + 1$, donc $\gamma_0 = 1/0.51$\index{processus AR(1)}. L'autocorrélation au retard $k$ vaut $\phi^k$. Un choc décroît comme $\phi^h$, donc la demi-vie vaut $\ln 2/(-\ln 0.7)$.

La meilleure prévision de $x_{t + h}$ est $\phi^hx_t$. Elle a pour variance $\phi^{2h}\gamma_0$, donc la part de variance expliquée vaut $\phi^{2h}$. Une longue simulation confirme la variance, l'autocorrélation et le second $R^2$.

Lorsque $\phi \to 1$, la demi-vie et la variance explosent, et le processus se rapproche d'une marche aléatoire.

\iqwrong{Utiliser $1/(1 - \phi)$ pour la variance, ce qui donne $3.3$. Elle vaut $1/(1 - \phi^2)$.}

\iqfollow{Quelle est la variance de l'erreur de prévision à deux pas~?}{$1 + \phi^2 = 1.49$, ce qui égale $\gamma_0(1 - \phi^4)$. L'erreur à long horizon tend vers la variance inconditionnelle.}

\iqwhy{Une aisance rapide avec le modèle de référence. La formule du $R^2$ montre à quelle vitesse la prévisibilité s'estompe.}
\end{solution}

\begin{solution}{iq:iv1:time-series:3}
La variance vaut $1.25$, l'autocorrélation de retard 1 vaut $0.4$, et celle de retard 2 vaut 0. La valeur maximale possible est $\frac12$, et le coefficient 2 donne les mêmes autocorrélations.

La variance vaut $(1 + \theta^2)\sigma^2$. La covariance au retard 1 vaut $\theta\sigma^2$, car seul $\varepsilon_{t - 1}$ est commun, donc $\rho_1 = \theta/(1 + \theta^2) = 0.5/1.25 = 0.4$\index{processus MA(1)}. Au-delà du retard 1, il n'y a aucun recouvrement, donc les autocorrélations s'annulent. Une simulation confirme ce résultat.

$\theta/(1 + \theta^2)$ atteint son maximum en $\theta = 1$, avec la valeur $\frac12$. Ainsi, toute autocorrélation de retard 1 supérieure à $\frac12$, sans autocorrélation au-delà du retard 1, ne peut pas provenir d'un MA(1).

Remplacer $\theta$ par $1/\theta$ conserve le rapport. Pour $\theta = 2$~: $2/5 = 0.4$. La convention retient $|\theta| < 1$, la forme inversible\index{inversibilité}, dans laquelle le bruit peut être reconstruit à partir des données passées.

\iqwrong{Donner $\rho_1 = 0.5$, c'est-à-dire le coefficient lui-même. Cela vaut pour un AR(1), mais pas ici.}

\iqfollow{D'où vient un MA(1) négatif dans les rendements~?}{Du bruit présent dans les prix observés. Si $p = x + u$, les différences comportent un terme $u_t - u_{t - 1}$, ce qui donne une autocovariance négative de retard 1.}

\iqwhy{C'est l'image miroir de l'AR(1) et la base de l'identification des modèles à partir des corrélations.}
\end{solution}

\begin{solution}{iq:iv1:time-series:4}
La journée la plus récente a un poids de 6~\%. La demi-vie vaut $11.2$ jours, le nombre effectif de jours est d'environ 32, et l'estimation de la volatilité est multipliée par $1.56$.

Les poids valent $0.06 \times 0.94^k$, dont la somme vaut 1\index{pondération exponentielle}. La demi-vie vérifie $0.94^k = \frac12$, donc $k = 11.2$.

La taille d'échantillon effective vaut $1/\sum w_k^2 = (1 + 0.94)/(1 - 0.94) = 32.3$. L'âge moyen des données vaut $1/0.06 = 16.7$ jours.

La mise à jour est $\hat{\sigma}_{\text{nouvelle}}^2 = 0.94\hat{\sigma}^2 + 0.06r^2$. Avec $r = 5\hat{\sigma}$, cela donne $(0.94 + 1.5)\hat{\sigma}^2 = 2.44\hat{\sigma}^2$. La volatilité augmente de $\sqrt{2.44} = 1.56$. Une seule journée a donc un effet important, et celui-ci persiste~: au bout de 11 jours, la moitié de la variance ajoutée subsiste.

\iqwrong{Dire que la volatilité augmente de $5 \times 6~\%$, ce qui confond volatilité et variance.}

\iqfollow{Quel est le lien avec le GARCH~?}{C'est un GARCH(1,1) intégré, avec $\omega = 0$, $\alpha = 0.06$, $\beta = 0.94$. Il n'existe aucune variance de long terme vers laquelle revenir, de sorte que la prévision est plate à tous les horizons.}

\iqwhy{C'est l'estimation de risque par défaut de nombreux desks, et le saut qui suit un mouvement important est ce qu'un gestionnaire des risques doit expliquer.}
\end{solution}

\begin{solution}{iq:iv1:time-series:5}
Il est stationnaire, avec des racines complexes $0.6 \pm 0.374i$ de module $0.707$ et une période de $11.3$ pas. Les autocorrélations valent $0.8$ et $0.46$.

La stationnarité dépend des racines de $z^2 - 1.2z + 0.5 = 0$\index{processus AR(2)}, qui valent $0.6 \pm \sqrt{0.36 - 0.5}$. Leur module vaut $\sqrt{0.5} = 0.707 < 1$, donc le processus est stationnaire.

L'argument vaut $\arctan(0.374/0.6) = 0.558$ radian, donc le cycle a pour période $2\pi/0.558 = 11.3$.

Les équations de Yule--Walker\index{equations de Yule-Walker@équations de Yule-Walker} donnent $\rho_1 = \phi_1 + \phi_2\rho_1$, donc $\rho_1 = 1.2/1.5 = 0.8$. Puis $\rho_2 = 1.2(0.8) - 0.5 = 0.46$. Une longue simulation donne le même résultat.

Un choc a pour effets 1, $1.2$, $0.94$, $0.53$, \dots. Il croît d'abord, puis décroît, devient négatif au pas 5, et oscille en se réduisant d'un facteur $0.707$ à chaque pas. C'est une onde amortie. Un AR(1) à coefficient positif ne peut pas faire cela. Notons que le coefficient $1.2$ dépasse 1 sans danger, car ce qui compte, ce sont les racines.

\iqwrong{Dire que $\phi_1 > 1$ le rend explosif, ou ne vérifier que $\phi_1 + \phi_2 < 1$, ce qui laisse de côté les autres conditions.}

\iqfollow{Un cycle du modèle signifie-t-il un cycle exploitable~?}{Pas à lui seul. L'amortissement est fort, et la phase de l'onde dérive au gré des nouveaux chocs, de sorte que le calage du cycle se perd en quelques périodes.}

\iqwhy{Lire la dynamique à partir des racines est la compétence clé. Des réponses en forme de bosse apparaissent en macroéconomie et dans les données de flux d'ordres.}
\end{solution}

\begin{solution}{iq:iv1:time-series:6}
Un AR(1) avec $\phi \approx 0.6$~: les autocorrélations décroissent géométriquement, et les autocorrélations partielles s'annulent brutalement après le retard 1. Les bandes de confiance valent $\pm 2/\sqrt{500} = \pm 0.089$.

$0.6^2 = 0.36$, $0.6^3 = 0.22$ et $0.6^4 = 0.13$, ce qui correspond à la première liste. L'autocorrélation partielle\index{autocorrélation partielle} au retard $k$ est le dernier coefficient d'un ajustement AR($k$), et c'est la corrélation qui subsiste une fois contrôlés les retards intermédiaires. Après le retard 1, toutes se situent dans les bandes, donc un seul retard suffit.

Un MA(1) présente le cas inverse~: l'autocorrélation s'annule après le retard 1, et les autocorrélations partielles décroissent, en alternant de signe si le coefficient est positif. Un ARMA(1,1) a les deux fonctions qui décroissent, sans que ni l'une ni l'autre ne s'annule\index{identification de modèle}.

Les simulations confirment cette signature~: un coefficient de régression de retard 2 proche de 0 pour l'AR(1), et proche de $-0.19$ pour un MA(1) de coefficient $0.5$.

Ces bandes supposent un bruit blanc. Sous l'AR(1) ajusté, les autocorrélations empiriques aux retards plus élevés ont une dispersion plus large, de sorte que les bandes sous-estiment l'incertitude à ces retards.

\iqwrong{Ajuster un ordre élevé parce que les autocorrélations restent au-dessus de la bande pendant plusieurs retards. Une décroissance lente ne dit rien à elle seule de l'ordre, et ce sont les valeurs partielles qui le révèlent.}

\iqfollow{Combien de retards paraîtront significatifs par hasard parmi 20, en l'absence de toute structure réelle~?}{Environ un, au seuil de 5~\%. Tester plusieurs retards ensemble est le rôle de la statistique de Ljung--Box.}

\iqwhy{C'est le tableau d'identification classique, et on peut le demander sur des résultats réels en une minute.}
\end{solution}

\begin{solution}{iq:iv1:time-series:7}
$1.1$ et $1.173$, avec une limite de $(1 + \rho)/(1 - \rho) = 1.222$. La volatilité réelle à long horizon est d'environ $10.6~\%$ supérieure à la valeur donnée par la règle de la racine carrée du temps.

La variance d'une somme de $q$ termes donne\index{ratio de variance}
\[
\mathrm{VR}(q) = 1 + 2\sum_{k = 1}^{q - 1}\left(1 - \frac{k}{q}\right)\rho_k.
\]
Avec $\rho_k = 0.1^k$~: $\mathrm{VR}(2) = 1 + 0.1$, et $\mathrm{VR}(5) = 1 + 2(0.8 \times 0.1 + 0.6 \times 0.01 + 0.4 \times 0.001 + 0.2 \times 0.0001) = 1.173$. Une longue simulation confirme ce résultat.

Pour $q$ grand, la somme tend vers $1 + 2\rho/(1 - \rho) = 1.222$. La volatilité évolue comme la racine du rapport, $\sqrt{1.222} = 1.106$. Ainsi, une petite autocorrélation quotidienne ajoute 10.6~\% à la volatilité annuelle, et elle transformerait un ratio de Sharpe annuel mesuré de 1 en environ $0.9$.

Avec $\rho < 0$, le rapport passe sous 1, comme pour des rendements à retour à la moyenne, qui paraissent moins bons sur données quotidiennes que sur de longs horizons.

\iqwrong{Utiliser $\mathrm{VR}(q) = 1 + (q - 1)\rho$, ce qui ignore à la fois les pondérations et la décroissance de $\rho_k$.}

\iqfollow{Les rendements mensuels d'un fonds ont une autocorrélation de $0.3$. Qu'est-ce que cela vous suggère~?}{Des cotations figées ou un lissage, plutôt que du talent, et cela signifie que la volatilité réelle est supérieure d'un facteur d'environ $\sqrt{1.857} = 1.36$.}

\iqwhy{C'est le test le plus simple de la marche aléatoire, et il a un effet direct sur tout chiffre annualisé.}
\end{solution}

\begin{solution}{iq:iv1:time-series:8}
La variance de long terme vaut $4 \times 10^{-5}$, soit 10~\% de volatilité annuelle, la demi-vie est de $13.5$ jours, et la prévision vaut $2.8$ fois la variance de long terme, soit $16.8~\%$.

En prenant l'espérance de la récurrence\index{GARCH}, $V = \omega + (\alpha + \beta)V$, donc $V = \omega/(1 - \alpha - \beta) = 2 \times 10^{-6}/0.05$. L'annualisation donne $\sqrt{252V} = 10\%$.

La prévision conditionnelle vérifie $\mathbb{E}\sigma_{t + h}^2 - V = (\alpha + \beta)^h(\sigma_t^2 - V)$. Les chocs s'atténuent donc au rythme de $0.95$ par jour, avec une demi-vie de $\ln 0.5/\ln 0.95 = 13.5$.

L'excès d'aujourd'hui vaut $3V$. Au bout de 10 jours, il vaut $0.95^{10} \times 3V = 1.8V$, donc la prévision vaut $2.8V$, soit une volatilité annuelle de $16.8~\%$. Une simulation de Monte Carlo sur la récurrence donne $2.8$.

Si $\alpha + \beta = 1$, $V$ n'existe pas, et la prévision reste plate.

\iqwrong{Appliquer la décroissance à la volatilité et non à la variance. La récurrence n'est linéaire qu'en variance.}

\iqfollow{Qu'est-ce que cela implique pour la structure par terme de la volatilité prévue~?}{Elle augmente ou diminue régulièrement vers 10~\% lorsque l'horizon s'allonge. C'est la variance moyenne sur $h$ jours qui compte pour une volatilité à terme.}

\iqwhy{Le GARCH(1,1) est la référence pour la prévision de volatilité. Il faut savoir lire la persistance et le niveau de long terme directement sur les coefficients.}
\end{solution}

\begin{solution}{iq:iv1:time-series:9}
La volatilité quotidienne vaut $1.19~\%$. Elle est environ cinq fois plus efficace qu'un seul rendement au carré. L'échantillonnage discret la biaise à la baisse, et le rebond bid-ask la biaise à la hausse.

$\ln 1.02 = 0.0198$, donc l'estimation vaut $0.0198^2/(4\ln 2) = 1.41 \times 10^{-4}$, dont la racine vaut $1.19\%$\index{estimateur de Parkinson}.

La constante vient de $\mathbb{E}R^2 = 4\ln 2\,\sigma^2$ pour l'étendue $R$ d'un mouvement brownien sans dérive. Un rendement au carré a une variance relative égale à 2. L'estimateur de Parkinson a une variance relative d'environ $0.41$, ce qui donne un rapport de $4.9$. Dans une simulation à pas fins, le rapport se situe entre $4.3$ et $5.5$.

Sur des données réelles, il se dérègle de deux façons principales. D'abord, le vrai plus haut et le vrai plus bas tombent entre deux cotations, de sorte que l'étendue observée est trop petite. Avec 78 cours par jour, une simulation montre que l'estimation de la variance est environ 14~\% trop basse, soit environ 7~\% en volatilité. Ensuite, les cotations au cours acheteur et au cours vendeur élargissent l'étendue, ce qui la tire vers le haut, davantage pour les titres peu liquides. Les mouvements de nuit sont aussi complètement manqués, de sorte que l'estimateur ne décrit que les heures de cotation, et la dérive le gonfle légèrement.

\iqwrong{Considérer le gain de cinq fois comme gratuit. Il repose sur le modèle, et les deux biais jouent l'un contre l'autre de façons qui dépendent du titre.}

\iqfollow{Comment combineriez-vous l'étendue et les rendements~?}{On utilise des estimateurs de type Garman--Klass, qui exploitent les cours d'ouverture, de plus haut, de plus bas et de clôture, et on ajoute un terme clôture-ouverture pour la nuit.}

\iqwhy{Les estimateurs par étendue sont peu coûteux et utiles, et le recruteur veut le chiffre d'efficacité accompagné des raisons de ne pas s'y fier.}
\end{solution}

\begin{solution}{iq:iv1:time-series:10}
La moyenne de l'estimateur vaut $\sigma^2 + 2n\eta^2$. Le biais relatif est de $+39~\%$ pour 78 rendements et de $+195~\%$ pour 390. Le bruit donne aux rendements une autocovariance d'ordre 1 égale à $-\eta^2$.

Un rendement observé vaut $\Delta x + \Delta u$\index{bruit de microstructure}, avec $\Delta u = u_t - u_{t - 1}$ de variance $2\eta^2$. En sommant $n$ rendements au carré, on obtient $\sigma^2$ provenant du prix efficient et $2n\eta^2$ provenant du bruit. Le terme de bruit croît avec $n$, alors que le signal, lui, ne croît pas.

Avec $\sigma^2 = 10^{-4}$ et $\eta^2 = 2.5 \times 10^{-7}$~: $2 \times 78 \times 2.5 \times 10^{-7} = 3.9 \times 10^{-5}$, soit $39~\%$ de $\sigma^2$. Pour 390, elle est cinq fois plus grande. Une simulation avec ces valeurs concorde avec la formule à 1~\%.

Deux rendements consécutifs partagent $u_t$ avec des signes opposés, de sorte que l'autocovariance vaut $-\eta^2$ quelle que soit la fréquence. Un graphique de la variance réalisée en fonction de la fréquence d'échantillonnage, le signature plot\index{signature plot}, croît à mesure que les pas rétrécissent.

L'estimateur de Roll interprète le même effet comme une fourchette~: avec un rebond de demi-fourchette $s/2$, $\eta = s/2$, et $s = 2\sqrt{-\operatorname{Cov}}$. Ici, cela donne une fourchette de $0.1~\%$.

\iqwrong{Supposer que davantage de données donne toujours une meilleure estimation. Dans ce cadre, la fréquence la plus élevée donne la pire estimation de toutes si rien n'est fait contre le bruit.}

\iqfollow{Comment peut-on utiliser l'autocovariance pour corriger la variance réalisée~?}{On ajoute le double de la somme des autocovariances d'ordre 1 à la variance réalisée, ce qui supprime le biais $2n\eta^2$ dans ce modèle. C'est l'idée qui sous-tend les noyaux réalisés.}

\iqwhy{C'est la première chose qui se dérègle avec des données haute fréquence. Le signe de l'autocovariance est un diagnostic utile.}
\end{solution}

\begin{solution}{iq:iv1:time-series:11}
Sous une racine unitaire, le régresseur n'est pas stationnaire, de sorte que la statistique admet une limite non normale. Les valeurs critiques sont $-1.95$ et $-2.87$. On s'attend à $\hat{\rho} \approx 0.979$.

Les asymptotiques usuelles reposent sur le fait que $\frac{1}{T}\sum x_{t - 1}^2$ tend vers une constante. Sous une marche aléatoire, cette quantité croît comme $T$ et reste aléatoire. La limite est donc un rapport de fonctionnelles du mouvement brownien\index{test de Dickey-Fuller}, et elle est asymétrique à gauche. Une simulation de $40\,000$ trajectoires avec $T = 250$ donne les quantiles à 5~\% égaux à $-1.95$ sans constante et $-2.87$ avec constante.

L'estimation elle-même est biaisée vers le bas\index{racine unitaire}. Avec une constante, $\mathbb{E}\hat{\rho} \approx 1 - 5.3/T$, soit $0.979$ pour $T = 250$, ce qu'une simulation confirme. Ainsi, une marche aléatoire paraît présenter un retour à la moyenne, avec une demi-vie apparente d'environ 33 jours. Un test utilisant la table normale au seuil de $-1.645$ rejetterait la vraie racine unitaire environ 45~\% du temps avec constante.

Supprimer la constante ne résout pas le problème~: la valeur à 5~\% reste loin de la loi normale, et elle suppose une moyenne nulle, hypothèse rarement sûre.

\iqwrong{Comparer la statistique à $-1.645$. C'est la façon la plus courante de découvrir des paires qui n'existent pas.}

\iqfollow{Pourquoi le biais est-il tellement pire avec une constante~?}{La série centrée emprunte au futur, de sorte que le régresseur est corrélé au bruit ultérieur. Il s'agit d'un décalage d'environ $5.3/T$ plutôt que d'environ $1.8/T$.}

\iqwhy{Toute stratégie de retour à la moyenne repose sur ce test. Connaître le biais et les valeurs critiques évite l'erreur la plus courante dans la recherche sur les paires.}
\end{solution}

\begin{solution}{iq:iv1:time-series:12}
La statistique vaut $2.0$ sous indépendance, et $1.63$ avec l'AR(1). La sélection du meilleur parmi de nombreux modèles rend même $2.0$ peu convaincant.

Le test de Diebold--Mariano\index{test de Diebold-Mariano} est une statistique $t$ sur la différence moyenne des pertes. Sous indépendance, elle vaut $0.02/(0.2/\sqrt{400}) = 2.0$, soit une $p$-valeur bilatérale de $0.046$.

Avec autocorrélation, la variance de la moyenne est multipliée par le facteur de long terme $(1 + \rho)/(1 - \rho) = 1.5$. La statistique vaut donc $2.0/\sqrt{1.5} = 1.63$, soit une $p$-valeur de $0.10$. Des séries AR(1) simulées avec ces paramètres donnent la même erreur type pour la moyenne. Les erreurs de prévision à des horizons supérieurs à un pas se chevauchent, et elles présentent normalement ce type de dépendance. Il faut donc un estimateur robuste à l'hétéroscédasticité et à l'autocorrélation, avec un nombre de retards adapté.

Si un modèle est le meilleur parmi 20 candidats, alors la meilleure statistique nulle attendue vaut déjà environ $1.87$. Une valeur de $2.0$ constitue alors à peine une preuve, et il faut un Reality Check de type data snooping, ou un échantillon de validation mis de côté\index{data snooping}. Le test compare aussi les prévisions telles quelles~; si les modèles sont emboîtés, ou si le bruit d'estimation compte, la loi limite n'est pas standard.

\iqwrong{Utiliser l'erreur type naïve pour des prévisions qui se chevauchent. Une $p$-valeur qui franchit le seuil de 5~\% ne le franchit souvent plus après cette correction.}

\iqfollow{Quelle est la différence entre l'ajustement dans l'échantillon (in-sample) et l'ajustement hors échantillon (out-of-sample) ici~?}{Ce test évalue des prévisions hors échantillon. L'ajustement dans l'échantillon favorise toujours le modèle le plus grand, et il ne dirait rien ici.}

\iqwhy{Toute affirmation d'amélioration exige une erreur type, et la comparaison de prévisions est la façon dont on choisit les modèles. Ces deux corrections couvrent l'essentiel des cas réels.}
\end{solution}

\begin{solution}{iq:iv1:time-series:13}
$x$ et $y$ sont I(1), et $u$ est I(0). La demi-vie vaut $6.6$ pas, et la pente converge à la vitesse $1/T$, et non $1/\sqrt{T}$. La méthode d'Engle--Granger régresse, puis teste les résidus avec une table spéciale.

$y$ et $x$ errent, alors que $y - 2x$ ne dérive pas. C'est une cointégration\index{cointégration} de vecteur $(1, -2)$. La demi-vie de $u$ vaut $\ln 2/(-\ln 0.9) = 6.6$.

La pente est superconvergente. Puisque $\sum x^2$ croît comme $T^2$, l'erreur $\sum xu/\sum x^2$ décroît comme $1/T$. Dans une simulation, quadrupler $T$ réduit l'erreur sur la pente d'un facteur d'environ 4, contre 2 dans une régression stationnaire.

Première étape~: régresser $y$ sur $x$ et prendre les résidus. Deuxième étape~: tester les résidus pour une racine unitaire. La table usuelle échoue, car $\hat{\beta}$ est choisi pour rendre les résidus aussi petits que possible, de sorte qu'ils paraissent stationnaires même quand ils ne le sont pas\index{méthode d'Engle-Granger}. Sur deux marches aléatoires indépendantes avec $T = 500$, la valeur critique à 5~\% ressort proche de $-3.37$, loin de $-1.95$ ou $-1.65$.

Avec deux marches aléatoires indépendantes, la régression est fallacieuse~: la pente reste aléatoire, et la statistique $t$ rejette la nullité de la pente plus de 70~\% du temps au seuil de 5~\%. La cointégration est le cas où une telle régression a un sens réel, et le test sur les résidus les distingue.

\iqwrong{Tester si la corrélation des niveaux est élevée, ou tester avec la table de Dickey--Fuller usuelle. La corrélation des rendements ou des niveaux ne dit rien sur le fait qu'un écart revienne.}

\iqfollow{Que peut-on faire avec plusieurs actifs~?}{On utilise la méthode de Johansen, une approche en système qui estime combien de combinaisons stationnaires existent, dans la limite du nombre d'actifs moins un.}

\iqwhy{C'est la base théorique du pairs trading. Un candidat doit savoir en quoi elle diffère de la corrélation et de la régression fallacieuse.}
\end{solution}

\begin{solution}{iq:iv1:time-series:14}
La kurtosis vaut $3.77$, finie lorsque $(\alpha + \beta)^2 + 2\alpha^2 < 1$. L'autocorrélation d'ordre~1 des carrés vaut $0.179$, et décroît de $0.95$ par retard.

Posons $V = \mathbb{E}\sigma^2$, $s = \alpha + \beta$ et $m = \mathbb{E}\sigma^4$. Élevons la récurrence au carré et utilisons $\mathbb{E}r^4 = 3\mathbb{E}\sigma^4$ et $\mathbb{E}[r^2\sigma^2] = m$, qui découlent de $r = \sigma z$ avec $z$ indépendant de $\sigma$\index{GARCH}~:
\[
m = \omega^2 + 2\omega sV + (3\alpha^2 + 2\alpha\beta + \beta^2)m.
\]
Remarquons que $3\alpha^2 + 2\alpha\beta + \beta^2 = s^2 + 2\alpha^2$. Avec $\omega = V(1 - s)$, le terme constant vaut $V^2(1 - s)(1 + s)$. Donc
\[
m = \frac{V^2(1 - s^2)}{1 - s^2 - 2\alpha^2}, \qquad \kappa = \frac{3m}{V^2} = \frac{3(1 - s^2)}{1 - s^2 - 2\alpha^2}.
\]
Avec $s = 0.95$ et $\alpha = 0.1$~: $3 \times 0.0975/0.0775 = 3.77$. La condition de finitude est un dénominateur positif. Ici $0.9025 + 0.02 < 1$, elle est donc satisfaite. Pour $\alpha = 0$, $\kappa = 3$, comme il se doit pour une loi normale.

Les carrés suivent un ARMA(1,1) de coefficient autorégressif $s$, et la formule donnée vaut $0.1 \times 0.1925/0.1075 = 0.179$. Chaque retard supplémentaire multiplie ce terme par $0.95$. Une simulation de deux millions de pas donne ces valeurs, et une autocorrélation des rendements proche de 0.

Les rendements n'ont donc aucune mémoire linéaire, alors que leurs amplitudes en ont une\index{regroupement de volatilité}. La valeur $0.179$ paraît modeste, mais la décroissance lente produit une persistance totale substantielle.

\iqwrong{Supposer que les queues épaisses exigent des innovations à queues épaisses. Ici des innovations normales produisent un excès de kurtosis par le mélange des variances, bien que les données réelles en demandent en général davantage.}

\iqfollow{Quelle est la kurtosis si $\alpha$ est doublé, $\alpha + \beta$ restant inchangé~?}{$3 \times 0.0975/(0.0975 - 0.08) = 16.7$. On est proche de la frontière~: pour $\alpha \approx 0.221$, elle devient infinie. Le quatrième moment est donc très sensible à $\alpha$, et la kurtosis empirique est une cible instable.}

\iqwhy{Ce calcul vérifie que le candidat sait manipuler les moments d'une récurrence. Il explique deux faits stylisés à la fois.}
\end{solution}

\begin{solution}{iq:iv1:time-series:15}
Le $R^2$ vaut $(\kappa/3 - 1)/(\kappa - 1) = 0.093$, même pour une prévision parfaite. Un $R^2$ faible face aux rendements au carré en dit donc peu, et un meilleur proxy de la variance est la solution.

Écrivons $r^2 = \sigma^2z^2$, avec $\mathbb{E}z^2 = 1$ et $z$ indépendant de $\sigma^2$. La régression de $r^2$ sur $\sigma^2$ a alors une pente égale à 1, et $R^2 = \operatorname{Var}\sigma^2/\operatorname{Var}r^2$.

En utilisant $\mathbb{E}r^4 = 3\mathbb{E}\sigma^4$ et $\mathbb{E}r^4 = \kappa V^2$~: $\operatorname{Var}\sigma^2 = (\kappa/3 - 1)V^2$ et $\operatorname{Var}r^2 = (\kappa - 1)V^2$. Pour $\kappa = 3.77$~: $0.257/2.77 = 0.093$. Une simulation du GARCH lui-même donne ce chiffre lorsque la prévision est la vraie $\sigma^2$.

Ainsi, des $R^2$ de 5~\% à 10~\% sont ce que produiraient des modèles parfaits, et les critiques qui les citent mesurent le bruit du proxy, lequel est énorme pour une estimation à partir d'une seule observation\index{proxy de volatilité}. Un modèle imparfait fait moins bien, ce pourquoi la formule constitue une borne supérieure pour ce modèle.

La solution consiste en un proxy moins bruité~: la variance réalisée à partir des rendements intrajournaliers. Avec 78 sous-rendements normaux, le bruit $\chi^2_{78}/78$ a pour variance $2/78$, et la formule devient
\[
R^2 = \frac{\kappa/3 - 1}{\kappa/3 - 1 + (\kappa/3)(2/78)} = 0.89.
\]
Une simulation confirme ce résultat à $0.05$ près. Comparer aussi les modèles avec une perte robuste au bruit du proxy, soit l'erreur quadratique, soit la perte quasi-vraisemblance, et éviter les erreurs absolues sur la volatilité, qui peuvent favoriser le mauvais modèle\index{perte QLIKE}.

\iqwrong{Conclure que le GARCH ne prévoit pas, à partir d'un $R^2$ de quelques pour cent face aux rendements au carré. Le test n'avait pratiquement aucune puissance au départ.}

\iqfollow{Que donne $\kappa = 3$, et pourquoi~?}{$R^2 = 0$~: la variance est constante, il n'y a donc rien à prévoir. La formule donne la limite cohérente.}

\iqwhy{C'est une confusion bien connue dans la recherche sur la volatilité, avec une résolution en deux lignes. Elle montre si le candidat distingue la qualité d'un modèle de sa mesure.}
\end{solution}

\begin{solution}{iq:iv1:time-series:16}
$\rho_k = R^2\phi^k$, donc $\rho_1 = 0.009$ et $\rho_{10} = 0.0035$. Le rendement suit un ARMA(1,1) avec $\phi = 0.9$ et un coefficient de moyenne mobile d'environ $0.89$. Les autocorrélations sont trop petites pour être vues, et pourtant elles s'additionnent.

Le bruit étant blanc et indépendant\index{espérances de rendement persistantes}, $\operatorname{Cov}(r_t, r_{t - k}) = \operatorname{Cov}(\mu_t, \mu_{t - k}) = \phi^k\operatorname{Var}\mu$. En divisant par la variance de $r$, on obtient $R^2\phi^k$.

Une décroissance géométrique à partir du retard~1 est la signature d'un ARMA(1,1), et un AR(1) plus bruit blanc en est exactement un. Sous la forme $r_t - \phi r_{t - 1} = e_t - \theta e_{t - 1}$, l'égalité de l'autocorrélation d'ordre~1 donne $\theta \approx 0.891$, qui annule presque la racine en $0.9$. D'où l'impression d'un bruit blanc.

Pour détecter $\rho_1 = 0.009$ à $t = 2$, il faut $(2/0.009)^2 \approx 49\,000$ observations, soit environ 196 ans de données quotidiennes. Une simulation de quatre millions de pas la fait apparaître.

L'effet est néanmoins réel. La somme des autocorrélations donne un ratio de variance à long terme de $1 + 2R^2\phi/(1 - \phi) = 1.18$. Les rendements à long horizon sont donc plus prévisibles que ceux à court horizon~: les corrélations sont faibles, mais elles ont toutes le même signe.

Tout tient à la manière dont la prévisibilité est atteinte. Si $\mu_t$ peut être observé à travers un signal distinct, le $R^2$ vaut 1~\%. La meilleure prévision à partir des seuls rendements passés n'atteint qu'environ $0.04~\%$, comme le montrent à la fois la variance d'innovation de l'ARMA et un filtre de Kalman en régime stationnaire. L'historique des rendements est donc un mauvais moyen de trouver un tel signal.

\iqwrong{Conclure, à partir d'autocorrélations non significatives, que les rendements sont imprévisibles. Une faible autocorrélation est compatible à la fois avec l'absence de signal et avec un signal persistant noyé dans un bruit important.}

\iqfollow{Comment exploiter cela sans signal externe~?}{Le plus souvent, on ne le peut pas, à ce $R^2$. Les régressions à long horizon gagnent un peu, mais elles réutilisent surtout les mêmes données et appellent des corrections de chevauchement.}

\iqwhy{Ce point résout un paradoxe apparent de la finance empirique, et il montre pourquoi les signaux externes comptent davantage que l'historique des rendements.}
\end{solution}

\enlargethispage{-3\baselineskip}%
\begin{solution}{iq:iv1:time-series:17}
Les parts stationnaires valent $\frac23$ et $\frac13$, la volatilité vaut $19.1~\%$, la kurtosis vaut $6.17$, et l'autocorrélation des carrés décroît comme $0.85^k$.

Pour les parts, équilibrons les flux\index{changement de régime}~: $\frac23 \times 0.05 = \frac13 \times 0.1$.

La variance est un mélange~: $\frac23(0.01) + \frac13(0.09) = 0.0367$, d'où une volatilité de $19.1~\%$. Une part de deux tiers dans l'état calme laisse tout de même l'essentiel de la variance à l'état tendu.

Pour la kurtosis, la loi normale donne $\mathbb{E}r^4 = 3\sigma^4$ dans chaque état~:
\[
\kappa = \frac{3(\frac23 \times 0.0001 + \frac13 \times 0.0081)}{0.0367^2} = 6.17.
\]
Cette kurtosis provient du seul mélange, sans queues épaisses dans aucun des deux états.

La mémoire de la chaîne est sa seconde valeur propre, $\lambda = 0.95 + 0.9 - 1 = 0.85$\index{chaîne de Markov!valeur propre}. Conditionnellement à l'état, $r_t^2$ et $r_{t + k}^2$ sont indépendants, donc leur covariance est celle des variances conditionnelles à l'état, qui décroît comme $\lambda^k$. La demi-vie vaut $4.3$ jours. Le niveau est fixé par la part de $\operatorname{Var}r^2$ due à l'état~: $\rho_k = 0.204\lambda^k$.

Une simulation de quatre millions de pas redonne $19.1~\%$, une kurtosis voisine de $6.2$, et la décroissance en $0.85$. Lorsque les états se rapprochent, la kurtosis tend vers 3, et l'autocorrélation vers 0.

\iqwrong{Moyenner les volatilités, $\frac23(10) + \frac13(30) = 16.7~\%$. Ce sont les variances qui se mélangent, pas les volatilités.}

\iqfollow{En quoi ce modèle diffère-t-il d'un GARCH de même persistance~?}{La mémoire est celle d'un état non observé, et non celle des rendements passés en tant que tels. Un mouvement de grande ampleur ne relève la prévision que par ce qu'il révèle sur l'état, ce qui peut produire un saut brusque, plafonné à la variance tendue, et non une hausse proportionnelle à $r^2$. À partir d'une probabilité a priori de $0.05$, un rendement de 4 écarts-types dans l'état calme porte la probabilité de l'état tendu à $0.956$, et la prévision de variance du lendemain est multipliée par $4.5$. Filtrer l'état est un problème d'inférence.}

\iqwhy{Les modèles à régimes sont souvent évoqués sans chiffres. Ceux-ci proviennent d'une chaîne à deux états et d'un peu d'algèbre.}
\end{solution}

\begin{solution}{iq:iv1:time-series:18}
$\mathbb{E}\hat{\phi} \approx 0.9115$, donc la demi-vie rapportée vaut environ $7.5$ pas, contre une valeur vraie de $13.5$. Plus de la moitié du biais provient de l'estimation de la moyenne, et un jackknife en retire l'essentiel.

Le biais vaut $-(1 + 3 \times 0.95)/100 = -0.0385$\index{biais de l'AR(1)}, d'où une estimation moyenne de $0.9115$. Les demi-vies valent $\ln 0.5/\ln\phi$~: $13.5$ pour $0.95$, et $7.5$ pour $0.9115$. Une simulation de $200\,000$ échantillons donne une estimation moyenne de $0.906$ et une demi-vie médiane d'environ $7.9$, donc à cette longueur le biais réel est un peu plus grand que celui de la formule.

Il y a deux sources. Même avec une moyenne connue, $\hat{\phi}$ est biaisé d'environ $-2\phi/T$, car le choc $\varepsilon_t$ entre dans les régresseurs ultérieurs. Le centrage ajoute environ $-(1 + \phi)/T$~: la moyenne de l'échantillon contient les termes de bruit ultérieurs, de sorte que $x_{t - 1}$ centré évolue en sens contraire de $\varepsilon_t$, et l'estimation baisse. Ici les formules donnent $-0.019$ et $-0.0195$, et la simulation $-0.017$ et $-0.027$. Près de 1, la correction est la plus forte, et avec un $T$ petit elle peut faire paraître réverting n'importe quelle marche aléatoire.

Ainsi, un backtest qui sélectionne les spreads à demi-vie courte favorise précisément ceux qui portent le plus de biais, et les profits espérés sont surestimés.

Pour le corriger, un jackknife à demi-échantillon\index{jackknife} utilise $2\hat{\phi} - \frac12(\hat{\phi}_1 + \hat{\phi}_2)$, où les deux dernières estimations proviennent des deux moitiés. Comme le biais est d'ordre $1/T$, cette combinaison en élimine l'essentiel. Dans la simulation, la moyenne revient à $0.953$. Utiliser directement la formule, en résolvant $\hat{\phi} = \phi - (1 + 3\phi)/T$ en $\phi$, donne $0.945$. Des estimations fondées sur la simulation ou médianement sans biais sont également employées. Toutes ont une variance plus grande, ce qui constitue leur coût.

\iqwrong{Rapporter la demi-vie à partir de l'estimation brute. Une demi-vie est une fonction fortement non linéaire près de $\phi = 1$, si bien que de faibles biais ont de grands effets.}

\iqfollow{Résolvez la formule pour cet échantillon.}{$\phi = (\hat{\phi}T + 1)/(T - 3) = (0.9115 \times 100 + 1)/97 = 0.95$, ce qui retrouve la vraie valeur par construction lorsque l'estimation égale la moyenne de la formule.}

\iqwhy{Les estimations de demi-vie déterminent le capital qu'un livre de pairs trading porte. Le biais joue dans le sens qui flatte chaque candidat.}
\end{solution}

\begin{solution}{iq:iv1:time-series:19}
$T(\hat{\rho} - 1) \to \frac12(W_1^2 - 1)/\int_0^1W^2$, et la statistique $t$ tend vers $\frac12(W_1^2 - 1)/\sqrt{\int_0^1W^2}$. Le numérateur est négatif avec une probabilité de $0.683$. La valeur à 5~\% vaut $-1.95$, et la moyenne est d'environ $-1.78$.

Partons de $\hat{\rho} - 1 = \sum x_{t - 1}\varepsilon_t/\sum x_{t - 1}^2$, où $x_t = \sum_{i \le t}\varepsilon_i$.

En élevant au carré, $x_t^2 = x_{t - 1}^2 + 2x_{t - 1}\varepsilon_t + \varepsilon_t^2$. En sommant,
\[
\sum_tx_{t - 1}\varepsilon_t = \frac12\left(x_T^2 - \sum_t\varepsilon_t^2\right).
\]
Par le théorème de Donsker\index{théorème central limite fonctionnel}, $x_{\lfloor Ts\rfloor}/(\sigma\sqrt{T}) \to W_s$. Et $\sum\varepsilon_t^2/T \to \sigma^2$. Ainsi, le numérateur divisé par $\sigma^2T$ tend vers $\frac12(W_1^2 - 1)$. Le dénominateur divisé par $\sigma^2T^2$ tend vers $\int_0^1W^2$.

En faisant le rapport, les facteurs $T$ donnent $T(\hat{\rho} - 1)$, et $\sigma$ se simplifie. Pour la statistique $t$, l'erreur type vaut $s/\sqrt{\sum x_{t - 1}^2}$, avec $s \to \sigma$~; elle tend donc vers le numérateur divisé par $\sqrt{\int W^2}$\index{racine unitaire}.

Le numérateur vaut $\frac12(\chi_1^2 - 1)$, un chi-deux décalé, et il est négatif lorsque $|W_1| < 1$, événement de probabilité $0.683$. Ainsi, $\hat{\rho} < 1$ environ deux fois sur trois, même avec une vraie racine unitaire. C'est la source de l'asymétrie à gauche.

Une simulation sur $T = 500$ et une évaluation directe de la fonctionnelle sur une grille fine concordent~: la moyenne de $T(\hat{\rho} - 1)$ est d'environ $-1.78$, et le quantile de la statistique $t$ vaut $-1.95$. La valeur normale $-1.645$ rejetterait $9.5~\%$ du temps, et non 5~\%. Le dénominateur a pour moyenne $\frac12$, et il reste aléatoire à la limite, ce qui explique à lui seul que la normalité échoue.

\iqwrong{Appliquer un théorème central limite au seul numérateur. La somme n'est pas gaussienne~: c'est un chi-deux décalé. Chaque choc est indépendant de son régresseur, mais la taille des termes est fixée par $x_{t - 1}$, qui croît et reste aléatoire~; la variance conditionnelle de la somme reste donc aléatoire, même après normalisation.}

\iqfollow{Que devient une constante dans la limite~?}{On écrit la limite sous la forme $\int_0^1W\,dW/\int_0^1W^2$. Puis $W$ est remplacé par $W^\mu = W - \int_0^1W$ dans les deux intégrales, si bien que le numérateur devient $\frac12(W_1^2 - 1) - W_1\int_0^1W$. La loi se décale vers la gauche, et la valeur à 5~\% passe à environ $-2.86$.}

\iqwhy{C'est une utilisation compacte des fonctionnelles browniennes en économétrie, et c'est la preuve derrière le test le plus utilisé du domaine.}
\end{solution}

\begin{solution}{iq:iv1:time-series:20}
L'erreur vaut environ $2\sigma^4/n + 4n^2\eta^4$, avec le meilleur $n = (\sigma^4/(4\eta^4))^{1/3} \approx 34$. L'estimateur à deux échelles utilise toutes les cotations et n'a pas de biais du premier ordre.

Avec $n$ rendements, le biais vaut $2n\eta^2$, donc son carré vaut $4n^2\eta^4$\index{fréquence d'échantillonnage optimale}. La variance provient surtout du prix efficient, et pour une volatilité constante elle vaut $2\sigma^4/n$. Il existe aussi des termes de bruit plus petits, d'ordre $n\eta^4$ et $\sigma^2\eta^2$. En annulant la dérivée, on obtient $n^3 = \sigma^4/(4\eta^4) = 40\,000$, donc $n = 34$~: environ un rendement toutes les 11 minutes sur une journée de 6~h~30. Une recherche sur la courbe complète, incluant les termes plus petits, aboutit aussi à peu près à ce résultat.

L'estimateur parcimonieux écarte la plupart des données. La variance réalisée à deux échelles\index{variance réalisée à deux échelles} les conserve. On divise les $N$ cotations en $k$ sous-grilles décalées, chacune comptant environ $\bar{n} = (N - k + 1)/k$ rendements. Soit $\overline{\mathrm{RV}}$ la moyenne de leurs variances réalisées, dont l'espérance vaut $\sigma^2 + 2\bar{n}\eta^2$, et soit $\mathrm{RV}^{\text{all}}$ la variance réalisée utilisant toutes les cotations, d'espérance environ $2N\eta^2$. Alors
\[
\mathrm{TSRV} = \overline{\mathrm{RV}} - \frac{\bar{n}}{N}\mathrm{RV}^{\text{all}}
\]
a pour espérance $\sigma^2$, car l'estimation sur toutes les cotations est essentiellement une estimation du bruit.

La simulation comporte $23\,400$ cotations avec ces paramètres, et $k = 300$. L'estimation sur toutes les cotations a une moyenne d'environ 118 fois la vraie valeur. Le TSRV a une moyenne à moins de 3~\% de celle-ci. Sa racine de l'erreur quadratique moyenne se situe entre 10~\% et 30~\% de $\sigma^2$, et elle est inférieure à celle de l'estimateur parcimonieux à $n = 34$, comme doit l'être un bon estimateur qui utilise davantage de données.

Ces résultats supposent un bruit indépendant. Le bruit réel est autocorrélé et dépend du flux d'ordres, ce qui justifie en partie l'emploi de noyaux plus élaborés. Et le temps des transactions donne souvent de meilleurs résultats que le temps horloge.

\iqwrong{Choisir une grille de 5 minutes par habitude, sans tenir compte du bruit. C'est souvent raisonnable, mais c'est le résultat de ce calcul pour un certain niveau de bruit, et non une loi.}

\iqfollow{Que devient le meilleur $n$ si le bruit double~?}{La formule montre $n \propto \eta^{-4/3}$, si bien que $n$ diminue d'un facteur $2^{4/3} = 2.5$, pour atteindre environ 14 rendements par jour.}

\iqwhy{Un compromis net, suivi de l'idée d'utiliser l'estimateur bruité pour se corriger lui-même. Cela relie directement la pensée fondée sur un modèle à ce qu'on fait des données tick.}
\end{solution}
